\documentclass[11pt,a4paper]{article}

% ---------- Encodage et langue ----------
\usepackage[utf8]{inputenc}
\usepackage[T1]{fontenc}
% \usepackage{lmodern}  % décommenter si disponible (police vectorielle)
\usepackage[french]{babel}
\usepackage[a4paper,margin=2.3cm]{geometry}

% ---------- Mathématiques ----------
\usepackage{amsmath,amssymb,amsthm}
\usepackage{mathtools}
\usepackage{enumitem}
\usepackage[most]{tcolorbox}
\usepackage[colorlinks=true,linkcolor=blue!60!black,urlcolor=blue!60!black]{hyperref}

% ---------- Macros ----------
\newcommand{\R}{\mathbb{R}}
\newcommand{\N}{\mathbb{N}}
\newcommand{\Z}{\mathbb{Z}}
\newcommand{\Q}{\mathbb{Q}}
\newcommand{\C}{\mathbb{C}}
\newcommand{\K}{\mathbb{K}}
\newcommand{\Sun}{\mathbb{S}^1}
\newcommand{\interieur}[1]{\mathring{#1}}
\newcommand{\adh}[1]{\overline{#1}}
\newcommand{\Int}{\operatorname{int}}
\newcommand{\diam}{\operatorname{diam}}
\newcommand{\dist}{\operatorname{dist}}
\newcommand{\Vect}{\operatorname{Vect}}
\newcommand{\Bf}{B_f}
\newcommand{\norm}[1]{\left\|#1\right\|}
\newcommand{\abs}[1]{\left|#1\right|}
\newcommand{\opnorm}[1]{|\!|\!|#1|\!|\!|}
\newcommand{\Lc}{\mathcal{L}}
\newcommand{\Cz}{\mathcal{C}^0}
\newcommand{\Cinf}{\mathcal{C}^\infty}
\newcommand{\Cc}{\mathcal{C}^0_c}
\newcommand{\ie}{c'est-à-dire}
\newcommand{\ps}[2]{\left\langle #1,#2\right\rangle}
\DeclareMathOperator{\Ima}{Im}
\DeclareMathOperator{\Ker}{Ker}
\DeclareMathOperator{\supp}{supp}
\DeclareMathOperator{\sgn}{sgn}

% ---------- Environnements ----------
\newtcolorbox{exo}[1]{
  enhanced, breakable,
  colback=black!3, colframe=black!70,
  fonttitle=\bfseries, title={#1},
  boxrule=0.6pt, arc=1pt, left=6pt, right=6pt
}
\newtcolorbox{pointcle}{
  enhanced, breakable,
  colback=blue!4, colframe=blue!55!black,
  fonttitle=\bfseries, title={Point clé},
  boxrule=0.8pt, arc=2pt, left=6pt, right=6pt
}
\newtcolorbox{piege}{
  enhanced, breakable,
  colback=red!4, colframe=red!60!black,
  fonttitle=\bfseries, title={Piège / difficulté},
  boxrule=0.8pt, arc=2pt, left=6pt, right=6pt
}
\newtcolorbox{commentaire}{
  enhanced, breakable,
  colback=gray!8, colframe=gray!55!black,
  fonttitle=\bfseries, title={Commentaire},
  boxrule=0.6pt, arc=2pt, left=6pt, right=6pt
}
\newtcolorbox{methode}{
  enhanced, breakable,
  colback=green!4, colframe=green!40!black,
  fonttitle=\bfseries, title={Méthode},
  boxrule=0.8pt, arc=2pt, left=6pt, right=6pt
}

\theoremstyle{definition}
\newtheorem*{solution}{Solution}

\title{\textbf{4MA305 -- Bases d'analyse fonctionnelle}\\[2mm]
\Large TD n\textsuperscript{o}3 -- Espaces de Banach (énoncés)}
\author{}
\date{Année universitaire 2026--2027}

\begin{document}
\maketitle

\begin{commentaire}
Feuille d'exercices sur le chapitre~3 du polycopié. La \textbf{partie A} (Exercices~1 à~6) est complémentaire des exercices du poly (corrigés dans le TD0). La \textbf{partie B} reproduit la \textbf{feuille officielle} du TD n\textsuperscript{o}3 (Exercices officiels~1 à~9). La \textbf{partie C} (Exercices~7 à~9) prépare l'examen. Le corrigé détaillé est dans \texttt{corrige\_TD3.pdf} ; l'idée est de chercher les exercices en conditions d'examen avant de le lire. Difficulté indicative, partie A : 1, 2, 4 (classiques d'examen) ; 3, 6 (applications directes des théorèmes) ; 5 (synthèse). Partie B : 1, 3, 4 (cours) ; 2, 5, 7, 9 (à savoir refaire) ; 6, 8 (courts).
\end{commentaire}

\section*{Partie A -- Exercices complémentaires}

\begin{exo}{Exercice 1 -- Série de Neumann et équation de Volterra}
Soit $E$ un espace de Banach et $T\in\Lc(E)$.
\begin{enumerate}[label=(\alph*)]
  \item On suppose $\sum_{n\ge0}\opnorm{T^n}<\infty$ (c'est le cas si $\opnorm T<1$). Montrer que la série $\sum_nT^n$ converge dans $\Lc(E)$, que $\mathrm{Id}-T$ est inversible d'inverse $\sum_nT^n$, et que si $\opnorm T<1$ alors $\opnorm{(\mathrm{Id}-T)^{-1}}\le\frac1{1-\opnorm T}$.
  \item Soit $A\in\Lc(E)$ inversible (d'inverse continu) et $B\in\Lc(E)$ avec $\opnorm{A-B}<\frac1{\opnorm{A^{-1}}}$. Montrer que $B$ est inversible. Qu'en déduit-on sur l'ensemble des opérateurs inversibles ?
  \item Soit $k\in\Cz([0,1]^2;\R)$, $M:=\norm k_\infty$, et $Vf(x):=\int_0^xk(x,y)f(y)\,dy$ pour $f\in E:=\Cz([0,1];\R)$ muni de la norme uniforme. On admet que $V\in\Lc(E)$. Montrer que $|V^nf(x)|\le\frac{M^nx^n}{n!}\norm f_\infty$, et en déduire que pour toute $g\in E$ l'équation $f-Vf=g$ admet une unique solution $f\in E$.
  \item Résoudre par cette méthode $f(x)-\int_0^xf(y)\,dy=1$.
\end{enumerate}
\end{exo}

\begin{exo}{Exercice 2 -- Deux applications de Banach--Steinhaus}
\begin{enumerate}[label=(\alph*)]
  \item Soit $a=(a_n)_{n\in\N}$ une suite de scalaires telle que la série $\sum_na_nx_n$ converge pour toute suite $x\in c_0(\N)$. Montrer que $a\in\ell^1(\N)$.
  \item Soit $1<p<\infty$. Même question en supposant la convergence pour tout $x\in\ell^p(\N)$ : montrer que $a\in\ell^{p'}(\N)$. Que dire pour $p=1$ ?
  \item Soient $E$, $F$ deux espaces de Banach et $B:E\times F\to\K$ bilinéaire, \emph{séparément continue} : pour tout $x$, $B(x,\cdot)$ est continue, et pour tout $y$, $B(\cdot,y)$ est continue. Montrer qu'il existe $C$ tel que $|B(x,y)|\le C\norm x\norm y$.
  \item Montrer que (c) est faux sans complétude : considérer $B(P,Q):=\int_0^1PQ$ sur $\R[x]$ muni de la norme $\norm P_1:=\int_0^1|P|$.
\end{enumerate}
\end{exo}

\begin{exo}{Exercice 3 -- Sous-espaces supplémentaires et projections continues}
Soit $E$ un espace de Banach, $F$ et $G$ deux sous-espaces vectoriels tels que $E=F\oplus G$ (somme directe algébrique), et $P$ la projection sur $F$ parallèlement à $G$.
\begin{enumerate}[label=(\alph*)]
  \item Montrer que si $P$ est continue, alors $F$ et $G$ sont fermés.
  \item Réciproquement, on suppose $F$ et $G$ fermés. Montrer que $P$ est continue, en appliquant le Corollaire~3.4 à $\Phi:F\times G\to E$, $(f,g)\mapsto f+g$.
  \item Redémontrer (b) par le théorème du graphe fermé.
  \item Soit $F$ un sous-espace fermé d'un EVN $E$ et $a\notin F$. Montrer que $F\oplus\K a$ est fermé. En déduire que la somme d'un sous-espace fermé et d'un sous-espace de dimension finie est fermée.
\end{enumerate}
\end{exo}

\begin{exo}{Exercice 4 -- Opérateurs injectifs à image fermée}
Soient $E$, $F$ deux espaces de Banach et $T\in\Lc(E,F)$ injectif.
\begin{enumerate}[label=(\alph*)]
  \item Montrer que $\Ima T$ est fermée si et seulement s'il existe $c>0$ tel que $\norm{Tx}\ge c\norm x$ pour tout $x\in E$.
  \item Étudier les exemples : $T(u):=\big(\frac{u_n}{n+1}\big)_n$ sur $\ell^2(\N)$ ; l'injection de $\ell^1(\N)$ dans $\ell^2(\N)$ ; le décalage à droite $S$ de $\ell^2(\N)$.
  \item Montrer que si $\norm{Tx}\ge c\norm x$ pour tout $x$ et si $\Ima T$ est dense, alors $T$ est bijectif d'inverse continu.
  \item Montrer que si $\norm{Tx}\ge c\norm x$ pour tout $x$ et si $S\in\Lc(E,F)$ vérifie $\opnorm{S-T}<c$, alors $S$ est injectif et d'image fermée.
\end{enumerate}
\end{exo}

\begin{exo}{Exercice 5 -- Un sous-espace fermé de $\Cz$ formé de fonctions $\mathcal C^1$ est de dimension finie}
Soit $F$ un sous-espace vectoriel \emph{fermé} de $\big(\Cz([0,1];\R),\norm\cdot_\infty\big)$ dont tous les éléments sont de classe $\mathcal C^1$.
\begin{enumerate}[label=(\alph*)]
  \item Montrer que l'application $D:F\to\Cz([0,1])$, $f\mapsto f'$, a un graphe fermé. En déduire qu'il existe $C$ tel que $\norm{f'}_\infty\le C\norm f_\infty$ pour toute $f\in F$.
  \item Montrer que la boule unité fermée de $F$ est compacte.
  \item Conclure que $F$ est de dimension finie.
  \item L'hypothèse « fermé » est-elle nécessaire ?
\end{enumerate}
\end{exo}

\begin{exo}{Exercice 6 -- Séries de fonctions dans $L^p$ : Riesz--Fischer en action}
Soit $\Omega$ un ouvert de $\R^d$, $1\le p<\infty$, et $(f_n)$ une suite de $L^p(\Omega)$ telle que $M:=\sum_n\norm{f_n}_p<\infty$.
\begin{enumerate}[label=(\alph*)]
  \item Montrer que $\sum_n|f_n(x)|<\infty$ pour presque tout $x$, que $f:=\sum_nf_n$ (définie presque partout) appartient à $L^p(\Omega)$, et que la série converge vers $f$ dans $L^p(\Omega)$.
  \item En déduire que $L^p(\Omega)$ est complet, puis que toute suite convergente dans $L^p(\Omega)$ admet une sous-suite convergeant presque partout.
  \item Soit $(r_n)_{n\in\N}$ une énumération de $\Q\cap\,]0,1[$ et $f(x):=\sum_n2^{-n}|x-r_n|^{-1/2}$. Montrer que $f\in L^1(]0,1[)$, que $f$ est finie presque partout, mais que $f$ n'est de carré intégrable sur aucun intervalle ouvert non vide de $]0,1[$.
\end{enumerate}
\end{exo}

\section*{Partie B -- Feuille officielle du TD n\textsuperscript{o}3}

\subsection*{Complétude}

\begin{exo}{Exercice officiel 1 (Exercice 3.1 du poly) -- Caractérisation des Banach par les séries}
Montrer qu'un EVN dans lequel toute série absolument convergente est convergente est un espace de Banach.
\end{exo}

\begin{exo}{Exercice officiel 2 (Exercice 3.4 du poly) -- $\Cz(\Omega)$ est complet, $\Cc(\Omega)$ ne l'est pas}
Soit $\Omega$ un ouvert non vide de $\R^d$. Montrer que $\Cz(\Omega)$ muni de la distance $d_{\Cz(\Omega)}$ de l'Exemple~2.5 est complet (espace de Fréchet). Montrer par contre que $\Cc(\Omega)$ n'est complet ni pour la norme uniforme, ni pour $d_{\Cz(\Omega)}$.
\end{exo}

\subsection*{Équicontinuité et compacité}

\begin{exo}{Exercice officiel 3 -- Équicontinuité et uniforme équicontinuité sur un compact}
Soient $(X_1,d_1)$ et $(X_2,d_2)$ deux espaces métriques, $X_1$ compact, et $A\subset\Cz(X_1;X_2)$. Montrer l'équivalence entre :
\begin{enumerate}[label=\alph*)]
  \item $A$ est \emph{uniformément équicontinue} : $\forall\varepsilon>0,\ \exists\delta>0,\ \forall f\in A,\ \forall x,y\in X_1,\ d_1(x,y)<\delta\Rightarrow d_2(f(x),f(y))<\varepsilon$ ;
  \item $A$ est \emph{équicontinue} : $\forall x\in X_1,\ \forall\varepsilon>0,\ \exists\delta>0,\ \forall f\in A,\ \forall y\in X_1,\ d_1(x,y)<\delta\Rightarrow d_2(f(x),f(y))<\varepsilon$.
\end{enumerate}
\end{exo}

\begin{exo}{Exercice officiel 4 -- Bornitude et équicontinuité passent à l'adhérence}
\begin{enumerate}[label=\alph*)]
  \item Soit $(X,d)$ un espace métrique et $A\subset X$. Montrer que $A$ est borné si et seulement si $\adh A$ est borné.
  \item Soient $(X_1,d_1)$, $(X_2,d_2)$ métriques, $X_1$ compact, et $A\subset\Cz(X_1;X_2)$ (muni de la distance uniforme). Montrer que $A$ est équicontinue si et seulement si $\adh A$ l'est, et de même pour l'uniforme équicontinuité.
\end{enumerate}
\end{exo}

\begin{exo}{Exercice officiel 5 -- Compacts de $\Cz([0,1];\R)$}
Déterminer si les ensembles suivants sont compacts dans $(\Cz([0,1];\R),\norm\cdot_\infty)$ :
\begin{enumerate}[label=(\alph*)]
  \item $A=\{f\text{ continue} : \norm f_\infty\le1\}$ ;
  \item $A=\{f\text{ polynôme} : \norm f_\infty\le1\}$ ;
  \item $A_N=\{f\text{ polynôme de degré}\le N : \norm f_\infty\le1\}$ ;
  \item $\adh A$, où $A=\{f\text{ dérivable} : |f'|\le1\}$ ;
  \item $\adh A$, où $A=\{f\text{ dérivable} : f(1)=2,\ |f'|\le1\}$ ;
  \item $\adh A$, où $A=\{f : |f(0)|\le3\ \text{et}\ |f(x)-f(y)|\le5|x-y|^{1/3}\ \forall x\ne y\}$ ;
  \item $A=\{f : |f(x)-f(y)|\le4|x-y|^2\ \forall x\ne y\}$ ;
  \item $A=\{f : |f(1/2)|+\frac{|f(x)-f(y)|}{|x-y|^{3/2}}\le4\ \forall x\ne y\}$.
\end{enumerate}
\end{exo}

\subsection*{Convolution et opérateurs}

\begin{exo}{Exercice officiel 6 (Exercice 3.7 du poly) -- Pas d'élément neutre pour la convolution dans $L^1(\R^d)$}
Montrer que la convolution, comme loi interne sur $L^1(\R^d)$, n'admet pas d'élément neutre.
\end{exo}

\begin{exo}{Exercice officiel 7 -- Convolution $L^p\star L^{p'}$}
Soient $1\le p,q\le\infty$ avec $\frac1p+\frac1q=1$, $f\in L^p(\R^d)$ et $g\in L^q(\R^d)$.
\begin{enumerate}[label=(\alph*)]
  \item Montrer que $f\star g$ est bornée et uniformément continue.
  \item Montrer que si $1<p,q<\infty$, $f\star g$ tend vers $0$ à l'infini.
\end{enumerate}
\end{exo}

\begin{exo}{Exercice officiel 8 -- Une forme linéaire de norme $1$ qui n'atteint pas sa norme}
Soit $(a_n)_{n\in\N}$ une suite de $]0,1[$ tendant vers $0$. On définit sur $\ell^1(\N)$ la forme linéaire $\langle\ell_a,\mathbf x\rangle:=\sum_{n\in\N}(1-a_n)x_n$.
\begin{enumerate}[label=\alph*)]
  \item Montrer que $\norm{\ell_a}_{\Lc(\ell^1(\N),\R)}=1$.
  \item Montrer que pour tout $\mathbf x$ de la sphère unité de $\ell^1(\N)$, $\langle\ell_a,\mathbf x\rangle<1$.
\end{enumerate}
\end{exo}

\begin{exo}{Exercice officiel 9 -- Limite simple d'une suite d'opérateurs}
Soient $(X,\norm\cdot_X)$ et $(Y,\norm\cdot_Y)$ deux EVN sur $\C$, $(f_n)_{n\ge1}\subset\Lc(X,Y)$ une suite convergeant en tout point, et $f(x):=\lim_nf_n(x)$.
\begin{enumerate}[label=(\alph*)]
  \item Montrer que $f$ est linéaire.
  \item Montrer que $\sup_{\norm x_X=1}\norm{f(x)}_Y\le\liminf_n\norm{f_n}_{\Lc(X,Y)}$.
  \item Montrer par un contre-exemple que $f$ n'est pas forcément continue.
  \item Montrer que $f$ est continue si $X$ est complet, ou si $(f_n)$ possède une sous-suite bornée.
\end{enumerate}
\end{exo}

\section*{Partie C -- Vers l'examen : questions de cours, vrai ou faux}

\begin{exo}{Exercice 7 -- Questions de cours}
\emph{À traiter sans document : énoncé complet, puis démonstration.}
\begin{enumerate}[label=(\alph*)]
  \item Montrer qu'un EVN est de Banach si et seulement si toute série absolument convergente y converge (Proposition~3.1 et Exercice~3.1).
  \item Énoncer et démontrer le théorème de Riesz (Théorème~3.1), en passant par le Lemme~3.1.
  \item Montrer que $\Lc(E,F)$ est un espace de Banach dès que $F$ en est un (Proposition~3.4).
  \item Énoncer et démontrer le théorème de Banach--Steinhaus (Théorème~3.7).
  \item Énoncer le théorème de l'application ouverte (Théorème~3.8, admis ici). En déduire le Corollaire~3.4 (continuité automatique de l'inverse), puis énoncer et démontrer le théorème du graphe fermé (Remarque~3.8).
\end{enumerate}
\end{exo}

\begin{exo}{Exercice 8 -- Vrai ou faux ?}
Pour chaque affirmation, dire si elle est vraie ou fausse, et justifier par une preuve courte ou un contre-exemple. $E$ et $F$ désignent des EVN.
\begin{enumerate}[label=\arabic*.]
  \item $\Cz([0,1])$ muni de la norme $\norm f_1=\int_0^1|f|$ est un espace de Banach.
  \item Un espace de Banach de dimension infinie n'admet pas de base (algébrique) dénombrable.
  \item La boule unité fermée de $\Cz([0,1])$ est compacte pour la norme uniforme.
  \item Une partie bornée et uniformément équicontinue de $\Cz([0,1])$ est compacte.
  \item Si $(f_n)$ est une suite de $\mathcal C^1([0,1])$ avec $\norm{f_n}_\infty+\norm{f_n'}_\infty\le M$ pour tout $n$, alors $(f_n)$ admet une sous-suite uniformément convergente.
  \item Si $E$ est de Banach et si $(T_n)$ est une suite de $\Lc(E,F)$ convergeant simplement vers $T$, alors $\opnorm{T_n-T}\to0$.
  \item Une limite simple d'applications linéaires continues de $E$ dans $F$ est continue.
  \item Une bijection linéaire continue $E\to F$ a un inverse continu.
\end{enumerate}
\end{exo}

\begin{exo}{Exercice 9 -- Vrai ou faux ? (suite)}
Même consigne. $E$ et $F$ désignent des EVN, $(e_n)_{n\in\N}$ la suite canonique des espaces de suites.
\begin{enumerate}[label=\arabic*.]
  \item Dans un EVN, toute série absolument convergente est convergente.
  \item Si $F$ est un espace de Banach, alors $\Lc(E,F)$ est un espace de Banach, même si $E$ n'est pas complet.
  \item $\ell^1(\N)$ muni de la norme $\norm\cdot_\infty$ est un espace de Banach.
  \item $\mathcal C^1([0,1])$ muni de la norme $\norm\cdot_\infty$ est un espace de Banach.
  \item Dans un EVN de dimension infinie, tout compact est d'intérieur vide.
  \item Un espace de Banach de dimension infinie n'est pas réunion dénombrable de parties compactes.
  \item La suite $f_n:x\mapsto\sin(nx)$ admet une sous-suite uniformément convergente sur $[0,2\pi]$.
  \item Si $f_n\to f$ dans $L^1(]0,1[)$, alors $f_n\to f$ presque partout.
  \item Pour $f\in L^1(\R)$, l'application $a\mapsto\tau_af$ est uniformément continue de $\R$ dans $L^1(\R)$.
  \item Si $(\rho_n)$ est une approximation de l'unité issue d'une fonction test et $f\in L^\infty(\R)$, alors $\rho_n\star f\to f$ dans $L^\infty(\R)$.
  \item Soient $E$ un espace de Banach, $D$ une partie dense de $E$ et $(T_n)$ une suite de $\Lc(E,F)$ telle que $\sup_n\norm{T_nx}<\infty$ pour tout $x\in D$. Alors $\sup_n\opnorm{T_n}<\infty$.
  \item Soient $F$ un espace de Banach, $D$ une partie dense de $E$ et $(T_n)$ une suite \emph{bornée} de $\Lc(E,F)$ telle que $(T_nx)_n$ converge pour tout $x\in D$. Alors $(T_nx)_n$ converge pour tout $x\in E$.
  \item Une application linéaire continue et surjective entre deux espaces de Banach envoie les fermés sur des fermés.
  \item Si deux normes $N_1\le C\,N_2$ rendent toutes deux $E$ complet, elles sont équivalentes.
  \item Si $E$ et $F$ sont des espaces de Banach et $T\in\Lc(E,F)$ est injectif, alors $\Ima T$ est fermée dans $F$.
  \item Soient $E,F$ des espaces de Banach et $T:E\to F$ linéaire telle que : si $x_n\to x$ dans $E$ et $Tx_n\to y$ dans $F$, alors $y=Tx$. Alors $T$ est continue.
\end{enumerate}
\end{exo}

\end{document}
